\chapter*{Abstract}
\addcontentsline{toc}{chapter}{Introduction} % serve per inserire l'Indice nel Table of contents

\begin{abstract}
In this laboratory project, we investigate the reconstruction and analysis of top-antitop (\( t\bar{t} \)) events arising from a hypothetical heavy neutral resonance (\( Z' \)) in the ATLAS detector at the Large Hadron Collider (LHC). We study proton-proton collisions at \(\sqrt{s} = 13~\mathrm{TeV}\), focusing on the \( Z' \rightarrow t\bar{t} \) decay channel in the context of searches for physics beyond the Standard Model. The analysis classifies final states into dileptonic, lepton+jets, and all-hadronic topologies, and utilizes key kinematic observables such as transverse momentum, missing transverse energy, pseudorapidity, and angular separation \( \Delta R \). Emphasis is placed on the identification of high-\( p_T \) objects, boosted topologies, and b-tagged jets. Detector components and their respective pseudorapidity coverages are reviewed in connection to their role in resolving final-state particles. The results demonstrate how the reconstruction of invariant masses and jet substructure can enhance sensitivity to heavy resonances and differentiate potential BSM signals from Standard Model backgrounds.
\end{abstract}
